\section{Contexto del problema}

Saber qué planta es la que aparece en una fotografía es el primer paso de muchas actividades: registrar biodiversidad, vigilar especies en riesgo, enseñar botánica o controlar la calidad de un producto agrícola. Hoy esa identificación depende, en gran medida, de que una persona con conocimiento especializado revise cada imagen, y ese recurso no alcanza para la cantidad de fotografías que se generan.

El caso más claro es la \textbf{ciencia ciudadana}. La plataforma iNaturalist ha acumulado más de 300 millones de observaciones subidas por personas comunes \parencite{inaturalist}; si cada una tuviera que ser identificada a mano por un experto, la mayoría quedaría sin clasificar. Esto es relevante para la conservación, ya que el informe \textit{State of the World's Plants and Fungi} estima que el 39.4\,\% de las especies de plantas está amenazado de extinción \parencite{kew2020}, y no se puede monitorear una especie si no se sabe dónde aparece. Los \textbf{jardines botánicos} y los museos enfrentan el mismo problema a menor escala: tienen colecciones y visitantes que preguntan qué están viendo, pero poco personal para responder.

También existe una necesidad económica. En México se cultivaron alrededor de 46 tipos de flores en 2021, con un valor de producción de 9,590 millones de pesos; tan solo el clavel aportó 839 millones, y el Estado de México generó el 81.4\,\% del valor de las flores de temporada otoñal \parencite{siap2022}. En el \textbf{comercio de flores}, distinguir variedades de forma automática ayuda a clasificar producto, verificar envíos y llevar inventarios sin depender de inspección visual.

El problema de fondo es que la diversidad es enorme: se estima que hay más de 5,000 especies de mamíferos, 10,000 de aves y 30,000 de peces, pero más de 400,000 tipos de flores \parencite{kagglePetalsMetal}. Muchas de ellas son visualmente muy parecidas, por lo que distinguirlas requiere experiencia, e identificarlas manualmente es lento y no escala cuando se tienen miles de fotografías. Por eso resulta valioso contar con modelos que reconozcan el tipo de flor de forma automática a partir de una imagen.

Este proyecto parte de la competencia \textit{Petals to the Metal}, de Kaggle, cuyo reto es construir un modelo que identifique el tipo de flor en un conjunto de imágenes. La competencia se limita a 104 clases de flores y corresponde a un problema de clasificación multiclase de imágenes en el que el modelo debe captar detalles finos de forma, estructura y color para distinguir entre categorías cercanas.

El problema tiene varias dificultades que lo hacen relevante desde un enfoque de redes neuronales. Primero, las clases son heterogéneas: algunas contienen un solo subtipo de flor y otras agrupan muchas subvariedades. Segundo, las imágenes no se capturaron en el mismo entorno, lo que agrega variación de fondo, escala e iluminación. Por último, el número de imágenes por clase no es uniforme (sección~\ref{sec:balance}), por lo que, dada la métrica de evaluación, el modelo debe funcionar bien tanto en las clases con pocos ejemplos como en las que tienen muchos.
